%% parts/p07/functional-interpretations.tex --- Phase-4 P7 chapter 7.6
%% New content: N-a-053, N-a-054, N-a-055.
%% Rules remain in registry/calculi.  The only local proved result is the
%% cut-free LJ last-rule lemma; the interpretation and HA metatheorems are
%% cited, with the exact source-relative HA boundary stated in the body.

\chapter{Functional Interpretations and Constructive Properties}
\label{sec:p07-functional-interpretations}
\chaptermark{Functional Interpretations}

Functional interpretations attach computational witnesses to derivations,
but the witness language, target theory, and soundness theorem must be named
separately.  This chapter fixes one Dialectica translation and Kleene's number
realizability, then states the disjunction and numerical-existence properties
at their honest evidence levels.  The principal sources are
\cite{Godel1958,Troelstra1973,AvigadFeferman1998,Kleene1945,
Troelstra1998,TroelstraVanDalen1988}.

Heyting arithmetic \(\mathsf{HA}\) has the exact standing fixed in
Definition~\ref{def:p07-heyting-arithmetic}; this chapter does not introduce a
second presentation of the theory.

\section{System T and the Dialectica interpretation}
\label{sec:p07-dialectica}

\begin{hbdefinition}{def:p07-system-t}{Godel's System T@G\"odel's System T}
The finite types form the least set containing the ground type \(0\) and
closed under the right-associative type former
\((\alpha,\beta)\mapsto\alpha\to_{\mathrm{ty}}\beta\).  The typed terms form
the least family containing variables of every finite type, the constants
listed below, and the application and typed-lambda-abstraction constructions
of the explicitly typed lambda-calculus on
\hyperref[sheet:stlc]{\texttt{sheet:stlc}}, now with a ground type \(0\), a
constant \(0:0\), successor \(S:0\to_{\mathrm{ty}}0\), and, for every finite type
\(\alpha\), a recursor
\[
 R_\alpha:
 \alpha\to_{\mathrm{ty}}
 (0\to_{\mathrm{ty}}\alpha\to_{\mathrm{ty}}\alpha)
 \to_{\mathrm{ty}}0\to_{\mathrm{ty}}\alpha.
\]
The zero and successor computation rules for this displayed argument order
live only on \hyperref[sheet:system-t]{\texttt{sheet:system-t}}.  That sheet,
rather than this body entry, fixes the operational surface used below.
This is \emph{System T}: a typed term calculus of primitive recursion at all
finite types.  It differs from \(\mathsf{HA}\), which is a first-order theory
of numbers, and from the quantifier-free theory \(\mathsf T_{\mathrm{qf}}\)
below, whose
formulas assert equality of System-T terms.  The distinction is necessary
because the Dialectica theorem extracts System-T terms while proving their
matrix property in \(\mathsf T_{\mathrm{qf}}\).
\index{System T@System \(T\)}
\index{primitive recursive functional!finite type}
\end{hbdefinition}

\begin{hbdefinition}{def:p07-dialectica-target}{The quantifier-free target theory \(\mathsf T_{\mathrm{qf}}\)}
The target \(\mathsf T_{\mathrm{qf}}\) is the exact quantifier-free,
many-sorted finite-type
theory imported from \cite[Section~2.2]{AvigadFeferman1998}, with its System-T
term surface linked to \hyperref[sheet:system-t]{\texttt{sheet:system-t}}.  Its
terms are those of System~T (equivalently, the cited combinatory presentation);
its atomic formulas are equalities of same-typed terms; and it has the cited
defining equations, equality principles, classical propositional base,
same-type term substitution, and open-formula induction principle.
Higher-type equality is read exactly as in the selected source presentation.
The symbol therefore names that imported formal target theory, not a locally
underspecified abbreviation and not merely the strongly normalizing term
calculus of Definition~\ref{def:p07-system-t}.  This entry does not reproduce
the target's complete axiom and rule surface; consequently every theorem that
uses it below remains citation-only.  The distinction is used in
Theorems~\ref{thm:p07-dialectica-soundness} and
\ref{thm:p07-dialectica-consistency}.
\index{System T@System \(T\)!quantifier-free target theory}
\end{hbdefinition}

\begin{hbdefinition}{def:p07-dialectica-translation}{A fixed Dialectica translation}
For every arithmetic formula \(\varphi\), define a finite-type formula
\[
 \varphi^{\mathrm D}\defeq
 \exists\vec x\,\forall\vec y\,
 \varphi_{\mathrm D}(\vec x,\vec y),
\]
where the witness tuple \(\vec x\) and challenge tuple \(\vec y\) may be
empty and the matrix \(\varphi_{\mathrm D}\) is quantifier-free in
\(\mathsf T_{\mathrm{qf}}\).  Every component of those tuples has the finite
type determined by the recursive clause in which it occurs; bound variables
are chosen fresh and every displayed substitution is capture avoiding.  If
\(\varphi^{\mathrm D}=\exists\vec x\forall\vec y,\varphi_{\mathrm D}\) and
\(\psi^{\mathrm D}=\exists\vec u\forall\vec v,\psi_{\mathrm D}\), the clauses are
\[
\begin{aligned}
 \chi^{\mathrm D}
   &\defeq \chi &&(\chi\text{ atomic}),\\
 (\varphi\land\psi)^{\mathrm D}
   &\defeq \exists\vec x,\vec u\,\forall\vec y,\vec v\,
       (\varphi_{\mathrm D}(\vec x,\vec y)\land
        \psi_{\mathrm D}(\vec u,\vec v)),\\
 (\varphi\lor\psi)^{\mathrm D}
   &\defeq \exists z,\vec x,\vec u\,\forall\vec y,\vec v\,
       ((z=0\land\varphi_{\mathrm D}(\vec x,\vec y))\lor
        (z=1\land\psi_{\mathrm D}(\vec u,\vec v))),\\
 (\varphi\to\psi)^{\mathrm D}
   &\defeq \exists\vec f,\vec g\,\forall\vec x,\vec v\,
       (\varphi_{\mathrm D}(\vec x,\vec g(\vec x,\vec v))\to
        \psi_{\mathrm D}(\vec f(\vec x),\vec v)),\\
 (\forall z\,\varphi(z))^{\mathrm D}
   &\defeq \exists\vec f\,\forall z,\vec y\,
       \varphi_{\mathrm D}(\vec f(z),\vec y,z),\\
 (\exists z\,\varphi(z))^{\mathrm D}
   &\defeq \exists z,\vec x\,\forall\vec y\,
       \varphi_{\mathrm D}(\vec x,\vec y,z).
\end{aligned}
\]
Negation is implication to \(\bottom\), so its clause is derived rather than
added.  The tag \(z\) is necessary to distinguish the selected disjunct;
the functionals \(\vec f\) transform witnesses and \(\vec g\) transform a
purported counterexample to the conclusion into a challenge to the premise.
This is an inductively defined translation of formulas, not formula
equivalence in the source logic.  Its clauses are used verbatim by the cited
soundness theorem.
\index{Dialectica interpretation}
\index{functional interpretation!Dialectica}
\end{hbdefinition}

\begin{hbtheorem}{thm:p07-dialectica-soundness}{cited=AvigadFeferman1998}{Dialectica soundness for HA}
Let \(\varphi\) be a closed arithmetic sentence.  If
\(\derives_{\mathsf{HA}}\varphi\), then
one can effectively extract a tuple of closed System-T terms \(\vec t\) such that
the quantifier-free target theory proves
\[
  \derives_{\mathsf T_{\mathrm{qf}}}\varphi_{\mathrm D}(\vec t,\vec y).
\]
The challenge variables \(\vec y\) remain free, equivalently universally
closed at the metalevel.
\end{hbtheorem}

This is Theorem~2.4.1 of \cite{AvigadFeferman1998}.  It is an effective
metatheorem about HA derivations and the fixed target, not an internal HA
reflection principle and not a blanket interpretation of analysis.

\begin{hbcorollary}{thm:p07-dialectica-consistency}{cited=AvigadFeferman1998}{Relative consistency consequence}
In the weak arithmetical base used to formalize the cited syntactic
translation,
\[
  \mathsf{Con}(\mathsf T_{\mathrm{qf}})\to\mathsf{Con}(\mathsf{HA}).
\]
\end{hbcorollary}

The cited corollary uses the translation's fixed treatment of falsity.  This
sentence records the scope and provenance of the imported implication; it is
not a local derivation.  The result is not an internal
\(\mathsf{HA}\)-proof of \(\mathsf{Con}(\mathsf{HA})\), nor does it assert the
absolute consistency of either theory.

\begin{proofboundary}{def:p07-dialectica-proof-boundary}{What the Dialectica citation carries}
A local proof would require induction on every rule and arithmetic axiom of
the selected HA calculus, typed construction of each extracted functional,
the quantifier-free decidability functional used for contraction, and the
recursor construction that handles induction.  None is replaced by the six
translation clauses.  The two metatheorems therefore remain cited and have no
annex items.
\index{proof boundary!Dialectica interpretation}
\end{proofboundary}

\section{Kleene number realizability}
\label{sec:p07-number-realizability}

\begin{hbdefinition}{def:p07-number-realizability}{Kleene number realizability}
Fix an acceptable numbering of the unary partial recursive functions, a
primitive-recursive pairing function, and projections \(\pi_0,\pi_1\).  For
a closed arithmetic formula \(\varphi\), the phrase
``the number \(e\) realizes \(\varphi\)'' is defined by structural recursion using
Table~\ref{tab:p07-number-realizability}.  It names a metalinguistic relation
between a natural number and an arithmetic formula; it is not the forcing or
satisfaction relation.  For formulas with numerical parameters, the cited
extension requires one index to realize all closed numeral instances
uniformly; that parameter clause is not silently identified with any one row
of the closed-formula table.  Pairing, program application, and definedness
are used by the soundness and ECT entries below \cite{Troelstra1998}.
\index{realizability!Kleene number}
\index{number realizability}
\end{hbdefinition}

\begingroup\small
\begin{longtable}{@{}p{.19\textwidth}p{.72\textwidth}@{}}
\caption{Kleene number-realizability clauses.}
\label{tab:p07-number-realizability}\\
\toprule
Formula & A number \(e\) realizes it exactly when \\
\midrule
\endfirsthead
\toprule
Formula & A number \(e\) realizes it exactly when \\
\midrule
\endhead
Atomic \(\chi\) & \(\chi\) is true in the standard natural numbers; the value of
\(e\) carries no further information \\
\(\varphi\land\psi\) & \(\pi_0(e)\) realizes \(\varphi\) and
\(\pi_1(e)\) realizes \(\psi\) \\
\(\varphi\lor\psi\) & \(\pi_0(e)=0\) and \(\pi_1(e)\) realizes \(\varphi\), or
\(\pi_0(e)=1\) and \(\pi_1(e)\) realizes \(\psi\) \\
\(\varphi\to\psi\) & For every \(a\) realizing \(\varphi\), the indexed
computation with program number \(e\) on input \(a\) terminates and its output
realizes \(\psi\) \\
\(\forall x\,\varphi(x)\) & For every \(n\in\N\), the indexed computation
with program number \(e\) on input \(n\) terminates and its output realizes the
result of substituting the standard numeral denoting \(n\) for \(x\) in
\(\varphi\) \\
\(\exists x\,\varphi(x)\) & \(\pi_0(e)\) is the numerical witness and
\(\pi_1(e)\) realizes the result of substituting the standard numeral denoting
\(\pi_0(e)\) for \(x\) in \(\varphi\) \\
\(\bottom\) and \(\neg\varphi\) & No number realizes \(\bottom\); negation is
\(\varphi\to\bottom\), so \(e\) realizes \(\neg\varphi\) exactly when no
number realizes \(\varphi\), in which case the implication clause is vacuous \\
\bottomrule
\end{longtable}
\endgroup

\begin{hbtheorem}{thm:p07-number-realizability-soundness}{cited=Kleene1945}{Soundness of number realizability for HA}
There is an effective transformation which, from an
\(\mathsf{HA}\)-derivation of a closed arithmetic sentence \(\varphi\), produces a
natural number realizing \(\varphi\).
\end{hbtheorem}

The cited proof is an induction over the complete HA derivation system.  The
clause table alone is not that induction.

\begin{hbdefinition}{def:p07-almost-negative-ect}{Almost-negative formulas and \(\mathsf{ECT}_0\)}
An arithmetic formula is \emph{almost negative} when it contains no
disjunction and every existential quantifier occurs immediately in front of
an atomic formula.  This syntactic class is needed because unrestricted
extended Church's thesis is inconsistent.  Fix an arithmetic formula
\(\chi(e,x,y)\) representing the graph of the selected acceptable numbering:
it holds exactly when the indexed computation with program number \(e\) on
input \(x\) terminates with output \(y\).  For almost-negative
\(\varphi(x)\), an instance of \emph{extended Church's thesis}
\(\mathsf{ECT}_0\) is the universal closure of
\[
 \forall x\bigl(\varphi(x)\to\exists y\,\psi(x,y)\bigr)
 \to
 \exists e\,\forall x\bigl(\varphi(x)\to
   \exists y\,(\chi(e,x,y)\land\psi(x,y))\bigr).
\]
The graph formula is an eliminable local abbreviation for the usual
arithmetization by Kleene's computation predicate and output function; its
standing is an arithmetic formula, not a new function symbol.  Thus
\(\mathsf{ECT}_0\) is an object-language scheme after expansion, not the
informal Church--Turing thesis.  The almost-negative restriction is used in
the characterization theorem below.
\index{extended Church thesis@extended Church's thesis}
\index{almost-negative formula}
\end{hbdefinition}

\begin{hbtheorem}{thm:p07-ect-characterization}{cited=Troelstra1998}{Formal number-realizability characterization}
In the formalized partial-term presentation of \cite{Troelstra1998}, every
instance of \(\mathsf{ECT}_0\) is number-realizable.  Moreover, for every
closed arithmetic sentence \(\varphi\), derivability of \(\varphi\) in that
source's exactly displayed partial-term calculus after adjoining
\(\mathsf{ECT}_0\) is equivalent to the existence of a standard numeral whose
realizability of \(\varphi\) is derivable in the underlying partial-term
calculus.
\end{hbtheorem}

The partial-term calculus and the almost-negative restriction are
load-bearing.  Replacing that cited calculus by an undefined local theory name,
or erasing the restriction, would turn a formal characterization into an
unsupported unrestricted scheme.

\begin{hbdefinition}{def:p07-markov-principle}{Markov's principle for decidable arithmetic predicates}
For an arithmetic predicate \(\varphi(x)\), the relevant instance of
\emph{Markov's principle} is the universal closure of
\[
 \forall x\,(\varphi(x)\lor\neg\varphi(x))
 \to
 \bigl(\neg\neg\exists x\,\varphi(x)\to\exists x\,\varphi(x)\bigr).
\]
The displayed decidability premise is part of the scheme.  This is an
object-language principle over arithmetic, distinct from the classical
metatheoretic step used to analyze realizers.
\index{Markov principle@Markov's principle}
\end{hbdefinition}

\begin{hbtheorem}{thm:p07-markov-realizability-boundary}{cited=TroelstraVanDalen1988}{Markov principle: realizability and independence boundary}
Ordinary Kleene number realizability validates the instances of
Markov's principle in Definition~\ref{def:p07-markov-principle}, under the
classical metatheory used for that semantics, while \(\mathsf{HA}\) does not
derive the full scheme.
\end{hbtheorem}

Consequently ordinary number realizability is not a countermodel proving
Markov's principle independent of HA.  The nonderivability half requires the
separate Kripke/model argument cited in
\cite{TroelstraVanDalen1988}; modified realizability or another named model
must be specified before it is used as an independence method.  Realizability
topoi are outside this row; see \cite{VanOosten2008} only for further reading.

\begingroup\small
\begin{longtable}{@{}p{.25\textwidth}p{.31\textwidth}p{.35\textwidth}@{}}
\caption{Number-realizability applications and their honest conclusions.}
\label{tab:p07-realizability-applications}\\
\toprule
Input & Supported conclusion & Unsupported inference \\
\midrule
\endfirsthead
\toprule
Input & Supported conclusion & Unsupported inference \\
\midrule
\endhead
HA derivation & Effective number-realizer extraction,
Theorem~\ref{thm:p07-number-realizability-soundness} & Truth of every
classical principle \\
Formal realizability plus almost-negative-domain \(\mathsf{ECT}_0\) & The
characterization in Theorem~\ref{thm:p07-ect-characterization} & Unrestricted
ECT, or the informal Church--Turing thesis as an axiom \\
Ordinary number realizability and Markov's principle & Realizability validates
the displayed MP scheme & A realizability countermodel to MP; a different
named construction is required \\
Realizability of \(\bottom\) & No number realizes \(\bottom\), giving the
corresponding relative consistency reading & An internal HA proof of its own
consistency \\
\bottomrule
\end{longtable}
\endgroup

\section{Disjunction and numerical existence}
\label{sec:p07-constructive-properties}

\begin{hblemma}{thm:p07-lj-disjunction-last-rule}{proved=annex:p7.functional.1}{Cut-free LJ disjunction extraction}
If there is a derivation on
\hyperref[sheet:lj]{\texttt{sheet:lj}} of the empty-antecedent sequent
\(\Rightarrow\varphi\lor\psi\) containing neither Cut nor Mix, then that sheet
has a derivation of \(\Rightarrow\varphi\) or a derivation of
\(\Rightarrow\psi\).
\end{hblemma}

The lemma is only the canonical last-rule step.  Turning an arbitrary
intuitionistic derivation into such a derivation additionally requires the
complete cut-elimination theorem for the exact calculus.

\begin{hbtheorem}{thm:p07-ipc-disjunction-property}{cited=TroelstraVanDalen1988}{Disjunction property for intuitionistic propositional logic}
For propositional formulas \(\varphi\) and \(\psi\), if
\(\derives_{\mathsf{LJ}}\varphi\lor\psi\), then
\(\derives_{\mathsf{LJ}}\varphi\) or
\(\derives_{\mathsf{LJ}}\psi\), where derivability is in the propositional
fragment of \hyperref[sheet:lj]{\texttt{sheet:lj}} with empty antecedent.
\end{hbtheorem}

\begin{hbtheorem}{thm:p07-ha-disjunction-property}{cited=TroelstraVanDalen1988}{Disjunction property for HA}
For closed arithmetic sentences \(\varphi\) and \(\psi\), if
\(\derives_{\mathsf{HA}}\varphi\lor\psi\), then
\(\derives_{\mathsf{HA}}\varphi\) or
\(\derives_{\mathsf{HA}}\psi\).
\end{hbtheorem}

\begin{hbtheorem}{thm:p07-ha-numerical-existence}{cited=TroelstraVanDalen1988}{Numerical existence property for HA}
Let \(\varphi(x)\) be an arithmetic formula whose only possible free variable is
the displayed number variable.  If
\(\derives_{\mathsf{HA}}\exists x\,\varphi(x)\), then for some \(n\in\N\),
the result of substituting the standard numeral denoting \(n\) for \(x\) in
\(\varphi\) is derivable in \(\mathsf{HA}\).
\end{hbtheorem}

These are metatheorems about the existence of derivations, not
object-language elimination rules.  In particular, extracting a number that
realizes a disjunct does not by itself produce an HA derivation of that
disjunct.  The frozen v1.1 normalization sketch covered neither arithmetic
induction nor equality and therefore cannot serve as a proof of the last two
theorems.

\begin{proofboundary}{def:p07-ha-extraction-source}{Source-relative status of the constructive-property proofs}
The three preceding metatheorems are imported from the fixed HA development in
\textcite[Vol.~I, Chapter~3, \S5.6, beginning at p.~139]{TroelstraVanDalen1988}.
That source proves the disjunction and numerical-existence properties for its
fully specified arithmetic formalism, including equality and induction.  This
handbook does not reproduce that formalism's proof-term grammar, equality
conversions, induction reductions, normalization proof, or erasure theorem;
consequently it does not invent a local proof-term calculus or label the three
results locally proved.

The cited extraction has the expected canonical-form shape: a normalized
closed proof of a disjunction determines one disjunct, and a normalized closed
proof of a numerical existential determines a numeral and a proof of its
instance.  This sentence explains the cited argument; it is not a proof from
the smaller inherited normalization theorem, which explicitly omits arithmetic
equality and induction.  For the purely propositional last-rule fact available
locally, use Lemma~\ref{thm:p07-lj-disjunction-last-rule} and its annex proof.
\index{disjunction property!source-relative proof}
\index{numerical existence property!source-relative proof}
\end{proofboundary}

The resulting treatments keep four neighboring topics distinct: the exact
System-T/Dialectica clauses and HA soundness theorem, the complete
number-realizability clause table and its Markov boundary, the matched
source-relative constructive-property results, and the separately sourced proof-mining
survey in the structural-extensions chapter.
